%!TEX root = ../main.tex

\begin{changemargin}{0.8cm}{0.8cm}

~\vfill{}

\section*{Abstrakt}
\vskip 0.5em
Sběratelské karetní hry (\ac{TCG}) představují globální trh, na kterém je indexace a třídění tisíců jednotlivých karet významnou logistickou výzvou. S rozšiřováním sbírek se ruční třídění stává neefektivním, zatímco komerční automatizovaná řešení jsou stále cenově nedostupná a spoléhají se na softwarové ekosystémy s uzavřeným zdrojovým kódem. Kromě toho tradiční textové metody identifikace, jako je optické rozpoznávání znaků (\ac{OCR}), v této oblasti často selhávají kvůli převaze lokalizovaných uměleckých fontů, různých jazyků a vizuálního šumu v okolí. Hlavní výzvou je nedostatek dostupných hardwarových a softwarových platforem s otevřeným zdrojovým kódem, které by byly schopné spolehlivé vizuální identifikace karet. Zde představujeme robotickou třídicí platformu s otevřeným zdrojovým kódem integrovanou s počítačovým zrakem využívajícím hluboké metrické učení, která identifikuje karty \ac{TCG} na základě celkové vizuální podobnosti. Aplikací Additive Angular Margin Loss (\ac{arcface}) na páteř \ac{ResNet}-50 systém zmapoval umělecká díla na kartách do 512dimenzionálního euklidovského prostoru, což umožnilo rychlé vyhledávání vektorů. Abychom překonali absenci označených datových sad z reálného světa, vyvinuli jsme proces generování syntetických dat, který míchal digitální šablony karet na stochastická pozadí a simuloval fyzické artefakty, jako je odlesk objektivu, šum senzoru a částečné zakrytí. Tato metodika účinně překlenula doménový posun mezi digitálními skeny a fyzickými objekty. Naše experimentální výsledky prokázaly 97,8\,\% přesnost vyhledávání Top-1 na neomezených fyzických fotografiích bez trénování na jakýchkoli reálných obrázcích. Předpokládáme, že tato open-source architektura poslouží jako základ pro dostupné robotické třídění a v budoucnu umožní automatizované generování rozsáhlých reálných datových sad \ac{TCG}.

\vskip 1em

{\bf Klíčová slova } \KlicovaSlova

\vskip 2.5cm

\end{changemargin}
